\documentclass[11pt]{article}
\usepackage[margin=1in]{geometry}
\usepackage{amsmath,amssymb}
\usepackage{graphicx}
\usepackage{verbatim}
\usepackage[hidelinks]{hyperref}
\graphicspath{{../}}

% ---------------------------------------------------------------------------
% EXECUTABLE RESEARCH PAPER -- methods/results log for the invertible
% normalizing flow project. Every methods item names the exact file that was
% run and every results item is \input directly from that step's outputs/
% folder, so the document always matches what was actually executed.
% Regenerate: run the step*.py scripts, then pdflatex this file (twice).
% This is the full-documentation version; prose compression and removal of
% file-level references happens later, closer to submission.
% ---------------------------------------------------------------------------

\title{Invertible Normalizing Flows for Generative Image Modeling:\\
Executable Methods and Results Log}
\author{Matt Tivnan\thanks{Repository:
\texttt{github.com/tivnanmatt/invertible\_normalizing\_flow\_20260831}}}
\date{\today}

\begin{document}
\maketitle

\begin{abstract}
This document is the executable methods-and-results log for the invertible
normalizing flow project. Each numbered step of the pipeline is one
\texttt{step*.py} script paired with one YAML config in \texttt{configs/} and
one output folder under \texttt{outputs/}; the tables and figures below are
included directly from those output folders, so the compiled PDF is an exact
record of what was run and what it produced. Step~1 (this version) builds
PyTorch datasets and dataloaders with deterministic train/val/test splits for
the standard generative-image-modeling benchmarks: MNIST, FashionMNIST,
CIFAR-10/100, SVHN, the MedMNIST 2D collection, CelebA, Imagenette, ImageNet,
AFHQ, FFHQ, CelebA-HQ, and LSUN bedroom.
\end{abstract}

\section{Execution environment}
\begin{itemize}
  \item Host: \texttt{axis03}; container: \texttt{recon-dev}
        (image \texttt{localhost:5000/recon:acr-20260819-v1};
        PyTorch 2.11.0+cu128, torchvision 0.26.0, CUDA 12.4, 2 GPUs).
  \item Repository mounted at \texttt{/dev\_ws/invertible\_normalizing\_flow\_20260831}
        inside the container (shared workspace
        \texttt{/home/staticct/recon-dev/workspace} on the host).
  \item Datasets live on the axis03 data drive at
        \texttt{/data/image\_benchmarks}, mounted into the container at the
        same path. This is a shared local repository of public benchmarks,
        reused across projects; nothing dataset-sized is stored in git.
  \item All steps are run as
        \texttt{docker exec recon-dev python /dev\_ws/invertible\_normalizing\_flow\_20260831/step<N>\_*.py}.
  \item Exact package versions, the resolved config, and the command line for
        the run shown in this PDF are recorded in
        \texttt{outputs/step1\_datasets/data/provenance.json}.
\end{itemize}

\section{Step 1: Datasets and dataloaders}

\subsection{Methods}
\begin{itemize}
  \item Code: \texttt{step1\_datasets.py} (repository root). Config:
        \texttt{configs/step1\_datasets.yml}. Outputs:
        \texttt{outputs/step1\_datasets/\{data,figures\}/}.
  \item Command executed for this document:
        \texttt{python step1\_datasets.py}
        (runs every dataset enabled in the config).
  \item The script exposes a library API used by later steps:
        \texttt{build\_dataset(name, cfg)} returns a bundle of three
        \texttt{torch.utils.data.Dataset} splits plus metadata, and
        \texttt{get\_dataloaders(bundle, cfg)} wraps them in
        \texttt{DataLoader}s (batch size, workers, pinning from the config).
  \item Auto-download datasets: MNIST, FashionMNIST, CIFAR-10, CIFAR-100,
        SVHN (torchvision); the MedMNIST 2D collection (\texttt{medmnist}
        package, official train/val/test); Imagenette (fast.ai 10-class
        ImageNet subset, used as a fast stand-in for ImageNet); AFHQ
        (StarGAN-v2 release, fetched from the official Dropbox link); CelebA
        (torchvision Google-Drive download, quota permitting).
  \item Manual-download datasets (access terms or hosting): ImageNet
        (ILSVRC2012), FFHQ, CelebA-HQ, LSUN bedroom. When absent, the script
        records them as \emph{unavailable} together with exact acquisition
        instructions (Section~\ref{sec:unavailable}) instead of failing.
  \item Train/val/test policy (deterministic, seeded from the config):
        official splits are used whenever they exist (MedMNIST, CelebA);
        datasets with only train/test get a validation set carved from train
        (10\%); datasets whose public split is train/val only (Imagenette,
        ImageNet, AFHQ) use the official val as \emph{test} and carve val
        from train; single-pool datasets (FFHQ, CelebA-HQ) are split
        80/10/10.
  \item Transforms: \texttt{ToTensor} to $[0,1]$ floats only; variable-size
        natural-image datasets additionally get
        \texttt{Resize}+\texttt{CenterCrop} to the per-dataset
        \texttt{image\_size} in the config. No normalization or dequantization
        is baked in here -- for flow training those belong to the model step
        so the data step stays lossless.
  \item Per dataset the run itself is the test: all three splits are
        instantiated, real batches are pulled through the dataloaders,
        per-channel mean/std are computed over up to 20 train batches, and
        three figures are written (sample grid per split, pixel-intensity
        histogram, class distribution where integer labels exist).
  \item Machine-readable outputs: one JSON per dataset,
        \texttt{summary.json}, \texttt{summary.csv},
        \texttt{provenance.json}, and the \LaTeX{} fragments
        (\texttt{summary\_table.tex}, \texttt{figures\_step1.tex},
        \texttt{unavailable\_step1.tex}) that are included verbatim below.
\end{itemize}

\subsection{Configuration (verbatim)}
{\small\verbatiminput{../configs/step1_datasets.yml}}

\subsection{Results: dataset summary}
\begin{center}
\InputIfFileExists{../outputs/step1_datasets/data/summary_table.tex}{}
  {\emph{outputs/step1\_datasets not generated yet -- run
   \texttt{python step1\_datasets.py}.}}
\end{center}
Split sizes are reported after the split policy above; \emph{image} is the
tensor shape $C\times H\times W$ delivered by the dataloaders; \emph{classes}
is the number of integer classes (``--'' for unconditional / multi-label
datasets).

\subsection{Results: unavailable datasets}
\label{sec:unavailable}
\InputIfFileExists{../outputs/step1_datasets/data/unavailable_step1.tex}{}
  {\emph{not generated yet.}}

\subsection{Results: figures}
\InputIfFileExists{../outputs/step1_datasets/data/figures_step1.tex}{}
  {\emph{outputs/step1\_datasets not generated yet.}}
\clearpage

\subsection{Analysis}
\begin{itemize}
  \item Coverage: the available benchmarks span $28\times28$ grayscale
        (MNIST, several MedMNIST members) through $64\times64$ RGB natural
        images and faces (CIFAR, Imagenette, AFHQ, CelebA), which is the
        standard ladder for likelihood-based generative models: density
        results in bits/dim on MNIST/CIFAR-10, then perceptual-quality
        results on faces/animals at higher resolution.
  \item The pixel-intensity histograms show the strongly peaked, saturated
        distributions typical of 8-bit images (large masses at 0 and 1,
        e.g.\ MNIST background, CelebA skin tones). For normalizing flows
        this directly motivates uniform dequantization and a logit
        transform before the first coupling layer -- to be implemented and
        ablated in the model step, not in the data step.
  \item Class-distribution figures document which datasets are usable for
        class-conditional modeling and how imbalanced they are; several
        MedMNIST members are markedly imbalanced, which matters for
        conditional likelihood evaluation.
  \item Per-channel means/stds recorded in each dataset's JSON give the
        constants for any input normalization later steps choose to apply.
  \item Datasets marked unavailable (typically ImageNet, FFHQ, CelebA-HQ,
        LSUN, and CelebA when the Google Drive quota blocks download) are
        access-restricted or manually hosted; the pipeline documents the
        exact acquisition path and picks them up automatically once placed
        under \texttt{/data/image\_benchmarks}.
\end{itemize}

\section{Step 2: TarFlow -- Transformer Autoregressive Flow}

\subsection{Methods}
\begin{itemize}
  \item Model: TarFlow from \emph{Normalizing Flows are Capable Generative
        Models} (Zhai et al., Apple, arXiv:2412.06329). A TarFlow is a stack
        of $T$ causal-ViT ``metablocks'', each an autoregressive NVP flow over
        image patches; the patch order is flipped between consecutive blocks.
        Training maximizes exact likelihood; sampling inverts the blocks
        patch-by-patch with KV caching.
  \item Code: the \emph{official} implementation
        (\texttt{github.com/apple/ml-tarflow}, Apple sample-code license) is
        used unmodified: \texttt{transformer\_flow.py} is imported directly
        from a clone at \texttt{/dev\_ws/ml-tarflow} (commit hash in
        \texttt{outputs/step2\_tarflow/data/provenance.json}). Our driver
        \texttt{step2\_tarflow.py} reimplements the surrounding training loop
        faithfully to the official \texttt{train.py} /
        \texttt{train\_local.ipynb} / \texttt{evaluate\_bpd.py}: AdamW
        ($\beta=(0.9,0.95)$, wd $10^{-4}$), cosine LR with one warmup epoch,
        inputs scaled to $[-1,1]$, gaussian-noise runs in bf16 autocast,
        uniform-dequantization (likelihood) runs in fp32, and the official
        bits/dim formula (Gaussian prior term with the $\log 128$ scaling
        correction, generalized from 3 channels to $C$).
  \item Data: datasets and deterministic splits come from step~1
        (\texttt{step1\_datasets.build\_dataset}); an adapter rescales to
        $[-1,1]$ and applies the official random horizontal flip where
        configured.
  \item Config: \texttt{configs/step2\_tarflow.yml}. Outputs:
        \texttt{outputs/step2\_tarflow/\{data,figures\}/}. Commands executed:
        \texttt{python step2\_tarflow.py --runs mnist\_toy} (GPU 0) and
        \texttt{python step2\_tarflow.py --runs cifar10\_uniform} (GPU 1),
        then \texttt{python step2\_tarflow.py --collect}.
  \item Run \textbf{mnist\_toy}: exact replication of Apple's released MNIST
        toy experiment (\texttt{train\_local.ipynb} defaults):
        class-conditional TarFlow $[4\hbox{-}128\hbox{-}4\hbox{-}4]$
        (patch-channels-blocks-layers), gaussian noise $\sigma=0.1$, batch
        256, lr $2\times10^{-3}$, 100 epochs. Deliverables: conditional
        sample grid (one row per digit) and the notebook's Bayes-rule
        classifier evaluation $\arg\min_y \frac12\mathbb{E}[z_y^2] -
        \log|\det|_y$ on the step-1 test split.
  \item Run \textbf{cifar10\_uniform}: unconditional CIFAR-10 likelihood run
        in the paper's likelihood configuration style: uniform
        dequantization, fp32, patch size 2 (as in their ImageNet~64
        2.99-bpd config), TarFlow $[2\hbox{-}256\hbox{-}8\hbox{-}4]$
        ($\sim$26M params), batch 128, lr $2\times10^{-4}$, 60 epochs on one
        RTX~4090. Validation bpd is tracked every 5 epochs; final test bpd
        uses the official formula on the step-1 test split.
  \item Scale caveat (stated up front): the paper's quantitative likelihood
        benchmark is \emph{unconditional ImageNet 64$\times$64}, where
        TarFlow $[2\hbox{-}768\hbox{-}8\hbox{-}8]$ ($\sim$470M params, 8
        $\times$ A100, 200 epochs) reaches \textbf{2.99 bpd}; the paper
        reports no MNIST or CIFAR-10 numbers. ImageNet is not yet staged
        (step~1), so this step validates the framework end-to-end at small
        scale and positions results against the classic CIFAR-10 flow
        literature (RealNVP 3.49, Glow 3.35, Flow++ 3.08 bpd).
  \item Sampling uses the plain reverse pass from Gaussian noise; the
        paper's score-based denoising step (Tweedie correction) is not
        applied here.
\end{itemize}

\subsection{Configuration (verbatim)}
{\small\verbatiminput{../configs/step2_tarflow.yml}}

\subsection{Results}
\begin{center}
\InputIfFileExists{../outputs/step2_tarflow/data/summary_table_step2.tex}{}
  {\emph{outputs/step2\_tarflow not generated yet -- run
   \texttt{python step2\_tarflow.py}.}}
\end{center}
Model column is [patch-channels-blocks-layers/block]. Full per-epoch metrics
are in \texttt{outputs/step2\_tarflow/data/*\_metrics.csv}.

\subsection{Results: figures}
\InputIfFileExists{../outputs/step2_tarflow/data/figures_step2.tex}{}
  {\emph{not generated yet.}}
\clearpage

\subsection{Analysis}
\begin{itemize}
  \item The mnist\_toy run reproduces Apple's released toy example: training
        loss decreases smoothly from the first epochs ($-0.70$ at epoch~1 to
        $\approx-1.62$ at epoch 100), the conditional sample grid shows clean,
        clearly class-consistent digits (one row per class), and the
        Bayes-rule classifier reaches \textbf{97.62\%} test accuracy --
        confirming the trained TarFlow is a working class-conditional density
        model, which is exactly the property the official notebook
        demonstrates. Residual background grain in the samples is the
        $\sigma=0.1$ training noise; the paper's denoising step (not applied
        here) is designed to remove it.
  \item The cifar10\_uniform run was \emph{interrupted at epoch 9/60}: GPU~1
        of the host failed under sustained load (fell off the PCIe bus,
        wedging the NVIDIA driver for new processes host-wide) and the run is
        queued to restart after a host reboot; the trainer now checkpoints
        optimizer state every epoch so future interruptions resume exactly. The
        partial trajectory is healthy and archived (metrics CSV + curves):
        train loss $-1.79 \to -3.18$ over 9 epochs and validation bpd
        \textbf{4.041} at epoch~5, the only bpd evaluation before the
        failure. For calibration: uniform-dequantization models start near
        8~bpd at random init, classic CIFAR-10 flows reach RealNVP 3.49 /
        Glow 3.35 / Flow++ 3.08 at convergence, and the paper's ImageNet-64
        state of the art is 2.99 bpd at $\sim$18$\times$ more parameters and
        $\sim$27$\times$ more GPU-epochs. Exact replication of the paper's
        headline number additionally requires staging ImageNet 64$\times$64
        (step-1 instructions).
  \item Framework verdict: the official TarFlow code trains stably out of
        the box under both noise regimes (bf16 gaussian, fp32 uniform), the
        likelihood pipeline (uniform dequantization + official bpd formula)
        produces sensible, steadily improving numbers, and inversion
        (sampling) works. The one incident so far was hardware, not the
        model.
\end{itemize}

\section{Step 3: Library of invertible modules with verification}

\subsection{Methods}
\begin{itemize}
  \item Code: \texttt{invlib.py} (the library) and
        \texttt{step3\_invertible\_modules.py} (the verification and
        benchmark harness). Config:
        \texttt{configs/step3\_invertible\_modules.yml}. Outputs:
        \texttt{outputs/step3\_invertible\_modules/\{data,figures\}/}.
  \item The library has two families with two contracts.
        \textbf{SeqTransform} implements TarFlow's permutation-slot signature
        \texttt{forward(x, dim, inverse)} and contains only ORTHOGONAL maps
        ($|\det|=1$): identity, flip, random permutation, orthonormal Haar
        wavelet, Walsh--Hadamard, DCT-II, discrete Hartley (real Fourier),
        fixed random rotations, and two learnable orthogonal
        parameterizations (products of Householder reflections; Cayley
        transform of a skew matrix), plus composition. Orthogonality is the
        exact condition under which the TarFlow metablock's likelihood
        accounting (no logdet term for the reordering) and the invariance of
        the $\mathcal{N}(0,I)$ prior are preserved, so these are drop-in
        replacements: the causal transformer then autoregresses over
        transform \emph{coefficients} -- arbitrary orthogonal features of the
        patch sequence -- instead of the raw patch order.
  \item \textbf{InvertibleModule} implements
        \texttt{forward(x)$\to$(y, logdet)} / \texttt{inverse(y)} with exact
        per-sample log-Jacobians: invertible diagonal (actnorm), Glow-style
        PLU dense linear and invertible $1{\times}1$ convolution
        (arXiv:1807.03039); invertible periodic (circular) convolutions
        decoupled in the frequency domain in 2D and 1D -- the ``invertible
        Fourier filter'' -- following Hoogeboom et al., arXiv:1901.11137;
        monotone elementwise cubic $x + ax^3$ ($a\ge0$) with closed-form
        Cardano inverse; general monotone odd polynomials (SOS-flow style,
        arXiv:1905.02325) inverted by bisection+Newton; invertible leaky
        ReLU; logit data transform; monotonic rational-quadratic splines
        (arXiv:1906.04032); Woodbury low-rank-plus-diagonal linear maps
        (Lu \& Huang, NeurIPS 2020); affine coupling (RealNVP-style); and
        \textbf{SpectralFloorLinear}, the eigenvalue-floored low-rank map
        proposed for this project: $A = U\,\mathrm{diag}(\lambda)U^\top +
        \lambda_0(I-UU^\top)$ with orthonormal $U\in\mathbb{R}^{D\times k}$
        and $\lambda_i = \lambda_0 + \mathrm{softplus}(\theta_i) \ge
        \lambda_0 > 0$ -- a full-rank invertible stand-in for a rank-$k$
        eigendecomposition in which every unretained eigenvalue is set to the
        floor $\lambda_0$ (necessarily $\le$ each retained one) instead of
        zero; the inverse and $\log\det = \sum_i\log\lambda_i +
        (D-k)\log\lambda_0$ are analytic.
  \item Verification per module (command: \texttt{python
        step3\_invertible\_modules.py}): (1) reconstruction
        $\max|x-f^{-1}(f(x))|$ in fp32 and fp64; (2) the module's logdet
        against $\log|\det|$ of the full autograd Jacobian (fp64, one
        sample); (3) $\|Q^\top Q - I\|_\infty$ for the orthogonal family;
        (4) a TarFlow integration test: each SeqTransform is installed in an
        actual \texttt{MetaBlock} (perturbed from identity) and the
        sequential reverse pass must invert the forward pass; (5) cost:
        median forward/inverse wall time and parameter count/bytes.
  \item The harness caught two real implementation bugs before any training
        used the library (a U/V swap in the Woodbury inverse, exposed by
        identical fp32/fp64 reconstruction error; and a determinant-sign
        over-restriction: the Hartley matrix has $\det=-1$, which is valid
        for flows since only $\log|\det|$ enters the likelihood).
\end{itemize}

\subsection{Configuration (verbatim)}
{\small\verbatiminput{../configs/step3_invertible_modules.yml}}

\subsection{Results: verification}
\begin{center}
{\small
\InputIfFileExists{../outputs/step3_invertible_modules/data/verification_table.tex}{}
  {\emph{not generated yet -- run \texttt{python step3\_invertible\_modules.py}.}}}
\end{center}

\begin{center}
\InputIfFileExists{../outputs/step3_invertible_modules/data/metablock_table.tex}{}
  {\emph{not generated yet.}}
\end{center}

\begin{figure}[htbp]
\centering
\includegraphics[width=0.85\textwidth]{outputs/step3_invertible_modules/figures/timing.png}
\caption{Median forward/inverse wall time per batch of 64 (CPU; the GPU rerun
uses the same harness). File:
\texttt{outputs/step3\_invertible\_modules/figures/timing.png}.}
\end{figure}

\begin{figure}[htbp]
\centering
\includegraphics[width=0.8\textwidth]{outputs/step3_invertible_modules/figures/orthogonal_matrices.png}
\caption{The orthogonal sequence transforms $Q$ ($64{\times}64$) available for
the TarFlow permutation slot. File:
\texttt{outputs/step3\_invertible\_modules/figures/orthogonal\_matrices.png}.}
\end{figure}

\begin{figure}[htbp]
\centering
\includegraphics[width=0.9\textwidth]{outputs/step3_invertible_modules/figures/mnist_coefficients.png}
\caption{An MNIST digit's patch sequence under each transform -- the
``features'' the autoregressive model would model. Haar concentrates the digit
into hierarchical detail coefficients; DCT/Hartley/Hadamard/random rotations
spread it into global features. File:
\texttt{outputs/step3\_invertible\_modules/figures/mnist\_coefficients.png}.}
\end{figure}
\clearpage

\subsection{Analysis}
\begin{itemize}
  \item All 26 modules PASS (24 original + the two 2D feature-domain
        tokenizers of step~5, verified here with roundtrip, norm
        preservation, DC-concentration and MetaBlock tests): fp32
        reconstruction errors are at roundoff
        ($10^{-7}$--$10^{-6}$), every analytic logdet matches the autograd
        Jacobian to $10^{-3}$ or (usually) far better in fp64, all orthogonal
        maps satisfy $\|Q^\top Q-I\|_\infty<10^{-5}$, and every SeqTransform
        passes the end-to-end MetaBlock forward$\to$reverse roundtrip at
        $\sim10^{-6}$ -- i.e.\ the whole orthogonal family is verified
        TarFlow-compatible before any training depends on it.
  \item Cost: fixed orthogonal transforms are parameter-free and as cheap as
        the baseline permutations at this scale (dense $64{\times}64$
        matmul); learnable orthogonal maps (Householder, Cayley) cost more
        per call because $Q$ is rebuilt each forward. Elementwise and
        FFT-based modules sit in between; inverse cost exceeds forward cost
        mainly for iterative inverses (monotone polynomial) and triangular
        solves (PLU).
\end{itemize}

\section{Step 4: Ablation -- what should TarFlow autoregress over?}

\subsection{Methods}
\begin{itemize}
  \item Code: \texttt{step4\_tarflow\_ablation.py}; config:
        \texttt{configs/step4\_tarflow\_ablation.yml}; outputs:
        \texttt{outputs/step4\_tarflow\_ablation/\{data,figures\}/}.
  \item Question: with everything else identical, which orthogonal feature
        basis in the permutation slot gives the best generative performance
        in a FIXED small number of epochs?
  \item Protocol: unconditional MNIST zero-padded to $32{\times}32$ (padding
        keeps 8-bit values valid for dequantization; $T=(32/4)^2=64$ is a
        power of two so Haar/Hadamard apply), TarFlow
        $[4\hbox{-}128\hbox{-}4\hbox{-}2]$, uniform dequantization in fp32,
        identical seed (model init and data order) for every variant. Odd
        blocks compose the transform with a flip so consecutive blocks
        autoregress in opposite coefficient orders, mirroring the official
        Identity/Flip alternation.
  \item Variants (one SeqTransform per metablock): baseline\_flip (official),
        identity\_only (control: no reordering), random\_perm, haar,
        hadamard, dct, hartley, rand\_ortho (per-block seeds), householder
        (learnable, 16 reflections), cayley (learnable).
  \item Tracked per variant: val bits/dim per epoch, final test bits/dim,
        wall time, transform parameter overhead, and (on GPU) peak memory;
        plus a sample grid from the reverse pass.
  \item Budget profiles: \texttt{full} (GPU: 120 epochs, all 54k train) and
        \texttt{reduced} (CPU: 6 epochs, 12k subset). The run reported here
        is the \texttt{full} GPU profile: all ten variants at 120 epochs on
        one RTX~4090, identical init, data order and budget. An earlier
        version of this document reported the \texttt{reduced} CPU profile
        (6 epochs, 12k subset), run while the host GPUs were unavailable;
        \S\ref{sec:step4analysis} records where the two disagree, because
        they disagree substantially.
\end{itemize}

\subsection{Configuration (verbatim)}
{\small\verbatiminput{../configs/step4_tarflow_ablation.yml}}

\subsection{Results}
\begin{center}
\InputIfFileExists{../outputs/step4_tarflow_ablation/data/ranking_table.tex}{}
  {\emph{not generated yet -- run \texttt{python step4\_tarflow\_ablation.py}.}}
\end{center}

\InputIfFileExists{../outputs/step4_tarflow_ablation/data/figures_step4.tex}{}
  {\emph{not generated yet.}}
\clearpage

\subsection{Analysis}
\label{sec:step4analysis}
\begin{itemize}
  \item The ranking table and overlay curve above are the study's primary
        readout; the run is a controlled comparison (same init, data order,
        budget), so ordering differences reflect the feature basis of the
        autoregression, not noise in the setup. The finite-epoch caveat
        applies by design: this measures which basis trains \emph{fastest},
        not which wins asymptotically.
  \item \textbf{The 6-epoch ranking did not survive the 120-epoch run, and
        the reversals are large.} Test bits/dim at 120 epochs:
        baseline\_flip 1.171, cayley 1.203, haar 1.608, random\_perm 2.157,
        householder 2.199, hadamard 2.780, dct 2.943, rand\_ortho 3.122,
        hartley 3.165, identity\_only 11.168. Against the reduced run:
        haar moves from 4.390 (reported as a failure, tied with the
        do-nothing control) to \emph{third place} at 1.608; the dense
        global bases move from 6.2--7.3 to 2.78--3.17; and identity\_only
        moves the other way, from 4.403 to 11.168. Every dense-basis
        variant was still improving at epoch 120 (e.g.\ haar gained 0.035
        bits/dim over its last 20 epochs versus the baseline's 0.020), so
        the 6-epoch numbers were measuring \emph{convergence speed}, not
        final quality. Any finite-budget ablation of this kind should be
        read with that confound in mind.
  \item A first pass was run WITHOUT gradient clipping (archived in
        \texttt{outputs/step4\_tarflow\_ablation/data\_noclip/}): ranking
        cayley 2.507 $<$ baseline\_flip 2.625 $\ll$ identity\_only 4.524 $<$
        random\_perm 4.773 $\ll$ hartley 6.23, dct 6.38, householder 6.56,
        hadamard 6.57, rand\_ortho 7.32, and haar diverged to NaN.
        \textbf{Correction:} an earlier version of this document attributed
        that instability to \emph{scale spread}, on the grounds that the DC
        coefficient is $\sqrt{T}\approx 8$ times the mean image value.
        Direct measurement on the data does not support this as the
        explanation for the \emph{ranking}. Pooled over 4000 test images,
        the Haar DC coefficient's standard deviation is only $1.5\times$ the
        median Haar coefficient's, and the Haar coefficient scales are in
        fact \emph{more} uniform across positions than raw patch scales.
        Scale spread remains a plausible account of the transient
        optimization blow-ups (which occur in every variant, including a
        $4.3\times10^{25}$ single-epoch spike in baseline\_flip at epoch~71
        that \texttt{grad\_clip} absorbed without lasting harm), but it does
        not explain which basis ends up ahead.
  \item \textbf{What does explain the fixed-basis ordering: destruction of
        exploitable sparsity, moderated by basis locality.} MNIST padded to
        $32\times32$ is largely deterministic background. Fitting the best
        \emph{diagonal} Gaussian in each basis (variance floored at the
        run's actual uniform-dequantization noise) gives the cost a
        factorized model would pay: raw patches $-2.464$ bits/dim with
        37.9\% of the 1024 dimensions near-constant; 1D Haar over the token
        sequence $-0.999$ with 17.6\%; 1D DCT/Hadamard/Hartley $+0.82$ to
        $+0.84$ with \emph{zero} near-constant dimensions. The mechanism is
        basis support: Haar basis vectors touch a median of 2 of the 64
        sequence slots, so a wavelet lying entirely inside the background
        still yields a constant coefficient, whereas every DCT and Hadamard
        basis vector spans all 64 slots and mixes background with digit.
        This proxy reproduces the measured ordering exactly, and predicts
        deficits (haar $+1.46$, dense $+3.28$ to $+3.31$) about a third to
        two-thirds larger than those observed --- as expected, since the
        proxy assumes a factorized model while the AR conditioning recovers
        some cross-coordinate structure. Note this is a penalty of
        optimization and inductive bias, not of information: orthogonal maps
        preserve differential entropy and $\mathcal{N}(0,I)$ is
        rotation-invariant, so all variants are equally expressive in
        principle.
  \item \textbf{And for the permutation variants: variance in
        autoregressive context coverage.} Counting, for each token, its
        total number of AR predecessors summed over the four metablocks, all
        three orderings give the same \emph{mean} (126) but different
        spreads: identity/flip alternation gives exactly 126 to every token
        (std 0) because the two orders are complementary by construction;
        four independent random permutations give std 36.3 with a
        worst-served token at 42; identity-only gives std 73.9 with token~0
        receiving \emph{zero} context in all four blocks. Since bits/dim is
        a sum over tokens and the cost of missing context is convex, the
        starved tokens dominate. That std ordering ($0 < 36.3 < 73.9$)
        reproduces the measured bpd ordering ($1.171 < 2.157 < 11.168$).
        Positional embeddings are \emph{not} the mechanism: the official
        implementation permutes the embedding table alongside the tokens, so
        each token keeps its own embedding under any ordering.
  \item \textbf{Experimental defect to fix before publication:}
        \texttt{random\_perm} is the only non-control variant not wrapped in
        the odd-block flip alternation that every other variant receives, so
        its deficit conflates random ordering with missing directional
        alternation. The clean version applies the \emph{same} permutation
        $\sigma$ on even blocks and $\mathrm{flip}\circ\sigma$ on odd
        blocks, which restores constant 126-token coverage while still
        destroying spatial locality, isolating the variable of interest.
  \item \textbf{cayley} -- a learnable rotation initialized near the
        identity ($Q\approx I$) -- remains the strongest alternative to the
        official ordering at 1.203 versus 1.171, and produces the cleanest
        sample grid in the study. The claim that it \emph{beats} the
        baseline does not survive: it won only in the unclipped 6-epoch run
        (2.507 vs 2.625) and finishes a close second at both 6 and 120
        epochs. The honest statement is that the official flip alternation
        is not beaten by any of the nine alternatives tested here, and that
        a near-identity-initialized learnable rotation matches it to within
        0.032 bits/dim. Householder ($k=16$, initialized far from identity)
        reaches 2.199 but costs $5\times$ the wall time (47.6 vs 9.4
        minutes) because its \texttt{matrix()} applies 16 reflections
        sequentially.
  \item \textbf{Likelihood and sample quality come apart.} Three variants
        produce entirely blank sample grids --- identity\_only (11.168),
        rand\_ortho (3.122) and householder (2.199) --- yet the latter two
        have competitive test likelihood, better than hartley (3.165), whose
        grid shows clear digits. A blank grid means the reverse pass
        saturated: every value clamped to the same bound. All three are
        unstructured dense rotations or the degenerate control, while the
        structured bases (Haar/DCT/Hadamard/Hartley) and the near-identity
        cayley all sample cleanly. TarFlow's reverse is sequential, so
        inversion error compounds across the chain; a good forward density
        does not imply a usable sampler. \emph{Rendering caveat:} the grids
        are written with \texttt{normalize=True} and no explicit
        \texttt{value\_range}, so each is contrast-stretched by its own
        min/max. Shape and texture are comparable across variants;
        brightness and contrast are not. Both should be fixed by passing
        \texttt{value\_range=(-1,1)} and by checkpointing weights so grids
        can be re-rendered without retraining.
  \item \textbf{The pixel-basis advantage is likely specific to MNIST.}
        Repeating the diagonal-Gaussian analysis on CIFAR-10, which has no
        deterministic background (zero near-constant dimensions in every
        basis), inverts every sign: relative to raw pixels, 1D Haar scores
        $-0.47$, full-image 2D Haar $-1.61$ and full-image 2D DCT $-2.17$
        bits/dim --- i.e.\ on natural images the pixel basis is the
        \emph{worst} of the four, not the best. Full-image 2D transforms
        also beat the 1D token-sequence versions on both datasets (MNIST
        $-0.75$, CIFAR $-1.15$), which supports the step-5 design. This is
        an analysis of the data, not a trained result: it predicts that the
        step-4 ranking will not transfer to CIFAR-10, and that prediction is
        untested until the ablation is actually run there.
  \item Designated follow-ups, in priority order: (i)~the same ablation on
        CIFAR-10, which the analysis above predicts will invert the ranking
        and is therefore the decisive test of whether any of this
        generalizes; (ii)~step~5's full-image 2D feature domains at the
        120-epoch budget, since they were only ever run at 6 CPU epochs ---
        the budget now known to penalize exactly these transforms;
        (iii)~learnable orthogonal transforms initialized AT the 2D Haar/DCT
        basis, combining the structure that helps the fixed bases with the
        adaptability that makes cayley competitive; (iv)~smoother wavelets
        (CDF~9/7, Daubechies-4) with better energy compaction than Haar's
        step functions; (v)~the corrected \texttt{random\_perm} variant
        described above; and (vi)~per-coefficient actnorm between blocks
        (as a proper flow layer with its logdet, not in the permutation
        slot) to remove the scale-spread confound entirely.
\end{itemize}

\section{Step 5: Feature-domain blocks -- autoregression over 2D
coefficients}

\subsection{Methods}
\begin{itemize}
  \item Code: \texttt{step5\_feature\_domains.py} (+
        \texttt{invlib.Haar2DFeatures}/\texttt{DCT2DFeatures}); config:
        \texttt{configs/step5\_feature\_domains.yml}; outputs:
        \texttt{outputs/step5\_feature\_domains/\{data,figures\}/}.
  \item Design correction relative to step~4: the step-4 transforms mixed
        token vectors across the sequence position only (a 1D transform on
        the patch index). Here the FEATURE EXTRACTION itself changes: each
        metablock folds its tokens back to the image, applies a full 2D
        orthogonal analysis (separable $HXH^\top$; multiresolution Haar or
        DCT-II), orders the coefficients coarse-to-fine (Haar: by
        max$(\mathrm{scale}_i,\mathrm{scale}_j)$; DCT: JPEG-style zigzag by
        $i{+}j$), chunks them into $T$ tokens, runs the causal AR update
        over those tokens, and synthesizes back to the image immediately
        after. Token~0 is the coarse approximation (for patch~4 on
        $32{\times}32$, exactly the 16 coarsest coefficients including DC);
        finer detail tokens condition on it -- a coarse-to-fine
        autoregression. Note the TarFlow metablock already applies the
        inverse of its reordering right after the causal transformer, so
        every block maps image~$\to$~image and heterogeneous domains stack.
  \item Both tokenizers are orthogonal on the flattened token space
        (fold/unfold and reordering are permutations; $H$ orthonormal), so
        $|\det|=1$, the likelihood accounting is unchanged, and they PASS
        the step-3 harness (roundtrip $\sim10^{-6}$, norm preservation,
        constant image $\Rightarrow$ 100\% of energy in token~0 with
        $\mathrm{DC}=N\cdot\mathrm{mean}$, MetaBlock reverse roundtrip).
  \item Variants (per-block domains, otherwise the step-4 protocol:
        identical init/data/budget, uniform dequantization, grad clip 1.0):
        \textbf{pixels\_flip} $[I,F,I,F]$ (official anchor);
        \textbf{haar2d\_alt} $[I,H,F,H{\circ}F]$ and \textbf{dct2d\_alt}
        $[I,D,F,D{\circ}F]$ (alternate pixel- and coefficient-domain
        blocks); \textbf{haar2d\_flip} / \textbf{dct2d\_flip} (all blocks in
        the coefficient domain, AR direction alternating);
        \textbf{haar2d\_pure\_dc} $[H,H,H,H]$ with NO flips.
  \item In haar2d\_pure\_dc the coarse token is causally first in every
        block and hence passes through the whole flow \emph{unchanged}. This
        yields the exact factorization $p(c) = p(c_0)\,p(c_{1:}\mid c_0)$:
        the flow is the conditional factor, and $p(c_0)$ (16 coefficients)
        is fit post-training as an independent per-dimension Gaussian KDE
        (Silverman bandwidth, 4k support points) that supports both
        log-density evaluation (entering the reported bits/dim exactly) and
        sampling (sampling draws $c_0$ from the KDE, Gaussian noise for all
        other coefficients, then runs the reverse pass, which leaves $c_0$
        untouched). The only approximation is independence across the 16
        dims of $c_0$; the factorization itself is exact.
\end{itemize}

\subsection{Configuration (verbatim)}
{\small\verbatiminput{../configs/step5_feature_domains.yml}}

\subsection{Results}
\begin{center}
\InputIfFileExists{../outputs/step5_feature_domains/data/ranking_table.tex}{}
  {\emph{not generated yet -- run \texttt{python step5\_feature\_domains.py}.}}
\end{center}

\InputIfFileExists{../outputs/step5_feature_domains/data/figures_step5.tex}{}
  {\emph{not generated yet.}}
\clearpage

\subsection{Analysis}
\label{sec:step5analysis}
\begin{itemize}
  \item \textbf{Budget warning --- read before comparing to step~4.} Every
        step-5 number in this section comes from the \texttt{reduced} CPU
        profile (6 epochs, 12k subset), and every step-4 number quoted
        \emph{within this section} is the corresponding step-4 reduced-profile
        number, not the 120-epoch results tabulated in
        \S\ref{sec:step4analysis}. The two are not comparable: at 120 epochs
        step~4's haar improves from 4.390 to 1.608 and its dense bases from
        6.2--7.3 to 2.78--3.17, so the step-5 conclusions below --- which
        rest on those reduced-profile contrasts --- are provisional and may
        reverse in the same way. Rerunning step~5 at the 120-epoch budget is
        the top-priority follow-up; the transforms it studies are precisely
        the slow-converging kind that the short budget penalizes.
  \item Protocol sanity: pixels\_flip reproduces the step-4 baseline
        \emph{at the matched reduced profile} to the fourth decimal (2.4736
        vs 2.4736 val bpd) -- the two steps share the protocol exactly, so
        numbers are directly comparable across both ablations whenever the
        budget is matched.
  \item Headline: true 2D coarse-to-fine feature extraction closes most of
        the gap that step-4's 1D sequence mixing opened.
        \textbf{Alternating pixel- and coefficient-domain blocks is nearly
        competitive with the official ordering} (dct2d\_alt 2.686,
        haar2d\_alt 2.823, vs baseline 2.474) -- compare step-4's
        all-coefficient 1D variants at 6.2--6.6. The heterogeneous stack
        gives the model complementary views: local raster AR in the pixel
        blocks, global coarse-to-fine AR in the coefficient blocks.
  \item All-coefficient stacks: haar2d\_flip 4.003 improves markedly on the
        step-4 1D seq-haar (4.390), while dct2d\_flip (6.463) does not
        improve on 1D dct (6.384) -- multiresolution locality (Haar
        coefficients have compact support) survives the domain change,
        global frequencies (DCT) do not. At this budget, pixel-domain blocks
        appear to be doing the heavy lifting wherever they are present.
  \item haar2d\_pure\_dc (4.342) sits between haar2d\_flip and the step-4
        1D haar: removing direction alternation costs a little, but the
        variant delivers what it was designed for -- an exact, interpretable
        factorization $p(c_0)p(\mathrm{rest}\mid c_0)$ whose reported bpd
        honestly includes the KDE term for the invariant coarse token, and
        whose samples draw global brightness/structure from the data-fitted
        KDE. Its sample grid is also the only one at this budget showing
        digit-like structure.
  \item Sampling note: for the best-bpd variants the sequential reverse pass
        at this (severely undertrained) stage emits rare huge values --
        sharper learned conditional scales amplify out-of-distribution
        prefixes -- which blanked the normalized sample grids; those grids
        are omitted from this version and the samplers now clip to $[-1,1]$
        like the official FID path. Sample quality at this budget is not a
        meaningful discriminator; bits/dim is the study metric, and the GPU
        \texttt{full}-profile rerun will regenerate all grids.
  \item GPU follow-ups: full profile for all six variants; haar2d\_alt with
        deeper stacks (more alternations); a joint (non-independent) model
        for $c_0$; and the same ablation on CIFAR-10 where global structure
        is richer.
\end{itemize}

\section{Reproducibility}
\begin{itemize}
  \item Global seed \texttt{20260831} (config) seeds torch, numpy, and every
        \texttt{random\_split}; splits are therefore identical across runs
        and machines.
  \item \texttt{outputs/} and dataset caches are gitignored; the committed
        code + config + this PDF are the record, and any output can be
        regenerated with the commands above.
  \item Provenance for the run in this PDF:
\end{itemize}
{\scriptsize
\IfFileExists{../outputs/step1_datasets/data/provenance.json}
  {\verbatiminput{../outputs/step1_datasets/data/provenance.json}}
  {\emph{provenance.json not generated yet.}}}

\begin{itemize}
  \item Provenance for step 2 (includes the pinned ml-tarflow commit):
\end{itemize}
{\scriptsize
\IfFileExists{../outputs/step2_tarflow/data/provenance.json}
  {\verbatiminput{../outputs/step2_tarflow/data/provenance.json}}
  {\emph{provenance.json not generated yet.}}}

\begin{itemize}
  \item Provenance for steps 3 and 4:
\end{itemize}
{\scriptsize
\IfFileExists{../outputs/step3_invertible_modules/data/provenance.json}
  {\verbatiminput{../outputs/step3_invertible_modules/data/provenance.json}}
  {\emph{provenance.json not generated yet.}}
\IfFileExists{../outputs/step4_tarflow_ablation/data/provenance.json}
  {\verbatiminput{../outputs/step4_tarflow_ablation/data/provenance.json}}
  {\emph{provenance.json not generated yet.}}}

\end{document}
